\begin{frame}{실제 숫자: 같은 환경, 계획이 있는가}
  \begin{itemize}
    \item 같은 격자 환경에서, \textbf{목표 상태(state 3, 오른쪽으로
      가는 행동)의 $Q$값이 0.5를 넘어서기까지}의 에피소드 수:
  \end{itemize}
  \vskip0.4em
  \small
  \begin{tabular}{@{}ll@{}}
    \toprule
    설정 & 0.5 초과까지 에피소드 \\
    \midrule
    \texttt{planning\_steps=0} (순수 Q-learning) & 253 \\
    \texttt{planning\_steps=10} & \textbf{9} $\leftarrow$ 매 실제
    스텝마다 모델로 10번 추가 학습 \\
    \bottomrule
  \end{tabular}
  \vskip1em
  \begin{block}{핵심}
    \kq{같은 \textbf{실제 경험}으로, 계획이 있는 쪽이 \textbf{훨씬
    더 많은 것}을 뽑아낸다.}
  \end{block}
\end{frame}
